\chapter{Second-Order Linear Differential Equations \& Oscillations}
\index{Second-Order ODEs}
\index{Harmonic Oscillator}
\index{Damped Oscillations}
\index{Overdamped}
\index{Critically Damped}
\index{Underdamped}
\index{Driven Oscillations}
\index{Resonance}
\index{Quality Factor}

\begin{chapterobjectives}
Upon mastering this chapter, students will be able to:
\begin{enumerate}
    \item Solve homogeneous linear second-order ODEs with constant coefficients using the characteristic auxiliary equation.
    \item Classify and physically interpret the three distinct damping regimes: overdamped, critically damped, and underdamped.
    \item Find particular solutions to inhomogeneous ODEs using undetermined coefficients and variation of parameters.
    \item Analyze driven mechanical and electrical oscillators, deriving the amplitude magnification and phase lag curves.
    \item Calculate the resonance frequency, resonance peak height, and bandwidth in terms of the Quality factor ($Q$).
    \item Compute and simulate damped and driven multi-frequency resonances numerically in Python.
\end{enumerate}
\end{chapterobjectives}

\section{Homogeneous Linear Second-Order ODEs with Constant Coefficients}
Second-order differential equations are central to physics because Newton's second law $\vec{F} = m\frac{d^2\vec{x}}{dt^2}$, wave equations $\nabla^2 \psi = \frac{1}{v^2}\frac{\partial^2 \psi}{\partial t^2}$, and Schrödinger's equation are all of second order.

The standard homogeneous equation with real constant coefficients is:
\begin{equation}
a\frac{d^2 x}{dt^2} + b\frac{dx}{dt} + c x = 0
\end{equation}
Substituting the trial solution $x(t) = e^{rt}$ yields the algebraic \textbf{characteristic equation}:
\begin{equation}
a r^2 + b r + c = 0 \implies r_{1,2} = \frac{-b \pm \sqrt{b^2 - 4ac}}{2a}
\end{equation}

\section{The Damped Harmonic Oscillator and Three Dynamic Regimes}
Consider a particle of mass $m$ attached to a spring of constant $k$ moving in a viscous damping medium with drag force $F_d = -b\dot{x}$. Newton's second law gives:
\begin{equation}
m\ddot{x} + b\dot{x} + kx = 0 \iff \ddot{x} + 2\gamma\dot{x} + \omega_0^2 x = 0
\end{equation}
where $\omega_0 \equiv \sqrt{k/m}$ is the **undamped natural frequency**, and $\gamma \equiv \frac{b}{2m}$ is the **damping coefficient**. The characteristic roots are:
\begin{equation}
r_{1,2} = -\gamma \pm \sqrt{\gamma^2 - \omega_0^2}
\end{equation}

Depending on the discriminant $\Delta = \gamma^2 - \omega_0^2$, the system exhibits three distinct physical regimes:
\begin{enumerate}
    \item \textbf{Overdamped ($\gamma > \omega_0$):} Roots are real, distinct, and negative. The mass returns exponentially to equilibrium without oscillating:
    \begin{equation}
    x(t) = C_1 e^{-(\gamma - \sqrt{\gamma^2 - \omega_0^2})t} + C_2 e^{-(\gamma + \sqrt{\gamma^2 - \omega_0^2})t}
    \end{equation}
    \item \textbf{Critically Damped ($\gamma = \omega_0$):} Repeated real root $r_1 = r_2 = -\gamma$. The system returns to equilibrium in the fastest possible time without overshoot:
    \begin{equation}
    x(t) = (C_1 + C_2 t)e^{-\gamma t}
    \end{equation}
    This is the ideal design for automobile shock absorbers and analog galvanometer meters.
    \item \textbf{Underdamped ($\gamma < \omega_0$):} Roots are complex conjugates $r_{1,2} = -\gamma \pm i\omega_d$, where $\omega_d \equiv \sqrt{\omega_0^2 - \gamma^2}$ is the \textbf{damped angular frequency}:
    \begin{equation}
    x(t) = A e^{-\gamma t}\cos(\omega_d t - \phi)
    \end{equation}
    The system oscillates with exponentially decaying amplitude envelopes $\pm A e^{-\gamma t}$.
\end{enumerate}

\begin{table}[H]
\centering
\small
\begin{tabularx}{\textwidth}{llccX}
\toprule
\textbf{Regime} & \textbf{Condition} & \textbf{Characteristic Roots} & \textbf{Oscillatory?} & \textbf{Engineering Application} \\
\midrule
Overdamped & $\gamma > \omega_0$ & Real, distinct: $r_1 \neq r_2 < 0$ & No & Heavy industrial door closers \\
Critically Damped & $\gamma = \omega_0$ & Real, repeated: $r = -\gamma$ & No & Vehicle suspension, seismographs \\
Underdamped & $\gamma < \omega_0$ & Complex: $-\gamma \pm i\omega_d$ & Yes & Quartz clocks, atomic force microscopy \\
\bottomrule
\end{tabularx}
\caption{Comparison of the three damping regimes in physical harmonic systems.}
\label{tab:ch09_damping_regimes}
\end{table}

\section{Driven Oscillations and Resonance}
When subjected to an external sinusoidal driving force $F(t) = F_0 \cos(\omega t)$, the equation of motion becomes:
\begin{equation}
\ddot{x} + 2\gamma\dot{x} + \omega_0^2 x = \frac{F_0}{m}\cos(\omega t)
\end{equation}
The steady-state particular solution is:
\begin{equation}
x_p(t) = A(\omega)\cos(\omega t - \delta)
\end{equation}
where the steady-state amplitude $A(\omega)$ and phase lag $\delta(\omega)$ are:
\begin{align}
A(\omega) &= \frac{F_0/m}{\sqrt{(\omega_0^2 - \omega^2)^2 + 4\gamma^2\omega^2}} \\
\delta(\omega) &= \arctan\left(\frac{2\gamma\omega}{\omega_0^2 - \omega^2}\right)
\end{align}

\begin{mathdefinition}[Resonance Frequency and Quality Factor]
The amplitude $A(\omega)$ reaches its absolute maximum at the \textbf{amplitude resonance frequency}:
\begin{equation}
\omega_r = \sqrt{\omega_0^2 - 2\gamma^2} \quad (\text{valid for } \gamma < \omega_0/\sqrt{2})
\end{equation}
The sharpness of the resonance curve is quantified by the dimensionless \textbf{Quality factor ($Q$)}:
\begin{equation}
Q \equiv \frac{\omega_0}{2\gamma} = \frac{\omega_0}{\Delta\omega_{\text{FWHM}}}
\end{equation}
where $\Delta\omega_{\text{FWHM}}$ is the full width at half-maximum of the energy resonance profile.
\end{mathdefinition}

\begin{pitfallbox}[Distinction Between Three Characteristic Frequencies]
Do not conflate the three distinct frequencies in damped systems:
\begin{itemize}
    \item Undamped natural frequency: $\omega_0 = \sqrt{k/m}$.
    \item Damped natural frequency: $\omega_d = \sqrt{\omega_0^2 - \gamma^2} < \omega_0$.
    \item Amplitude resonance frequency: $\omega_r = \sqrt{\omega_0^2 - 2\gamma^2} < \omega_d < \omega_0$.
\end{itemize}
Only in the weakly damped limit ($\gamma \ll \omega_0$ or $Q \gg 1$) do these three frequencies converge: $\omega_0 \approx \omega_d \approx \omega_r$.
\end{pitfallbox}

\section{Worked Examples and Physical Applications}

\begin{psctexample}[Driven Oscillator Amplitude and Phase Response]
An instrument package of mass $m = 2.0\text{ kg}$ is mounted on springs with effective constant $k = 800\text{ N/m}$ and damping $b = 4.0\text{ N}\cdot\text{s/m}$. A harmonic shaker applies driving force $F_0 = 10\text{ N}$ at frequency $\omega = 18\text{ rad/s}$.
Calculate: (a) undamped frequency $\omega_0$, (b) damping ratio $\gamma$, (c) quality factor $Q$, and (d) steady-state amplitude $A$ and phase lag $\delta$.

\textbf{\color{rbruNavy}Picture / Conceptualization:}
The mass is driven harmonically near its natural resonance. Damping limits the peak amplitude.

\textbf{\color{rbruNavy}Strategy / Governing Laws:}
Use the definitions $\omega_0 = \sqrt{k/m}$, $\gamma = b/(2m)$, and the steady-state response formulas for $A(\omega)$ and $\delta(\omega)$.

\textbf{\color{rbruNavy}Calculation / Analytical Solution:}
1. Natural frequency and damping parameter:
\begin{align}
\omega_0 &= \sqrt{\frac{800}{2.0}} = 20\text{ rad/s} \\
\gamma &= \frac{4.0}{2(2.0)} = 1.0\text{ s}^{-1}
\end{align}
2. Quality factor:
\begin{equation}
Q = \frac{\omega_0}{2\gamma} = \frac{20}{2(1.0)} = 10
\end{equation}
3. Steady-state amplitude at $\omega = 18\text{ rad/s}$:
\begin{align}
\omega_0^2 - \omega^2 &= 20^2 - 18^2 = 400 - 324 = 76\text{ (rad/s)}^2 \\
2\gamma\omega &= 2(1.0)(18) = 36\text{ (rad/s)}^2 \\
\text{Denominator} &= \sqrt{(76)^2 + (36)^2} = \sqrt{5776 + 1296} = \sqrt{7072} \approx 84.1\text{ (rad/s)}^2 \\
A(18) &= \frac{F_0 / m}{\text{Denominator}} = \frac{10 / 2.0}{84.1} = \frac{5.0}{84.1} \approx 0.0595\text{ m} = 5.95\text{ cm}
\end{align}
4. Phase lag:
\begin{equation}
\delta = \arctan\left(\frac{36}{76}\right) = \arctan(0.4737) \approx 0.442\text{ rad} = 25.3^\circ
\end{equation}

\textbf{\color{rbruNavy}Think / Physical Insights:}
Because $\omega < \omega_0$, the driving frequency is below resonance, and the response lags by less than $90^\circ$ ($\delta = 25.3^\circ$). At exact resonance ($\omega = \omega_0 = 20\text{ rad/s}$), the phase lag would be exactly $90^\circ$, and amplitude would reach $A_{\text{max}} \approx \frac{F_0/m}{2\gamma\omega_0} = \frac{5.0}{40} = 0.125\text{ m} = 12.5\text{ cm}$.
\qedexam
\end{psctexample}

\begin{pythonbox}[Driven Harmonic Resonance Curves in Python]
import numpy as np
import matplotlib.pyplot as plt

omega_0 = 20.0
F0_m = 5.0
omega = np.linspace(5, 35, 1000)

for Q in [2, 5, 10, 20]:
    gamma = omega_0 / (2 * Q)
    A = F0_m / np.sqrt((omega_0**2 - omega**2)**2 + 4 * gamma**2 * omega**2)
    print(f"Q = {Q:2d}: Peak Amplitude = {np.max(A)*100:.2f} cm")
\end{pythonbox}

\begin{psctexample}[Driven RLC Resonance and Q-Factor Analysis]
A series RLC circuit has $L = 10\;\text{mH}$, $C = 100\;\mu\text{F}$, $R = 2\;\Omega$, and is driven by $\mathcal{E}(t) = 12\cos(\omega t)\;\text{V}$. (i) Find the resonance frequency $\omega_0$, (ii) compute the Q-factor and bandwidth $\Delta\omega$, (iii) calculate the maximum current amplitude at resonance, and (iv) identify the analogous mechanical oscillator.

\textbf{Solution:}

\textbf{(i) Resonance frequency:}
\begin{equation}
\omega_0 = \frac{1}{\sqrt{LC}} = \frac{1}{\sqrt{10 \times 10^{-3} \times 100 \times 10^{-6}}} = \frac{1}{\sqrt{10^{-6}}} = 1000\;\text{rad/s}
\end{equation}
Corresponding to $f_0 = \omega_0/2\pi \approx 159.2\;\text{Hz}$.

\textbf{(ii) Q-factor and bandwidth:}
\begin{equation}
Q = \frac{\omega_0 L}{R} = \frac{1000 \times 0.010}{2} = \boxed{5}
\qquad
\Delta\omega = \frac{\omega_0}{Q} = \frac{1000}{5} = 200\;\text{rad/s}
\end{equation}
The $Q = 5$ indicates moderate selectivity: at frequencies $|\omega - \omega_0| > 100\;\text{rad/s}$, the response amplitude falls below $1/\sqrt{2}$ of peak.

\textbf{(iii) Maximum current at resonance} (where $X_L = X_C$ cancel):
\begin{equation}
\tilde{Z}(\omega_0) = R + i(\omega_0 L - 1/\omega_0 C) = R + 0 = 2\;\Omega
\implies \boxed{I_\text{max} = \frac{\mathcal{E}_0}{R} = \frac{12}{2} = 6\;\text{A}}
\end{equation}

\textbf{(iv) Mechanical analogy:}
The RLC equation $L\ddot{Q} + R\dot{Q} + Q/C = \mathcal{E}_0\cos\omega t$ maps exactly to $m\ddot{x} + b\dot{x} + kx = F_0\cos\omega t$ via: $L \leftrightarrow m$, $R \leftrightarrow b$, $1/C \leftrightarrow k$, $Q \leftrightarrow x$, $\mathcal{E}_0 \leftrightarrow F_0$.

\textbf{\color{rbruNavy}Think / Physical Insights:}
The Q-factor simultaneously quantifies: (1) energy stored relative to energy dissipated per cycle, $Q = 2\pi E/\Delta E_\text{cycle}$; (2) frequency selectivity (bandwidth); and (3) oscillation longevity after the drive is switched off ($\tau_\text{decay} = Q/\omega_0$). High-$Q$ resonators (quartz crystal, optical cavities) are fundamental to precise timekeeping and spectroscopy.
\qedexam
\end{psctexample}

\section{Summary of Key Formulas}
\begin{table}[H]
\centering
\small
\begin{tabularx}{\textwidth}{llX}
\toprule
\textbf{Oscillator Parameter} & \textbf{Formula} & \textbf{Physical Meaning} \\
\midrule
Natural Frequency & $\omega_0 = \sqrt{k/m}$ & Oscillation frequency without damping or driving \\
Damped Frequency & $\omega_d = \sqrt{\omega_0^2 - \gamma^2}$ & Actual oscillation rate under viscous resistance \\
Amplitude Resonance & $\omega_r = \sqrt{\omega_0^2 - 2\gamma^2}$ & Frequency of maximum steady-state displacement response \\
Quality Factor ($Q$) & $Q = \frac{\omega_0}{2\gamma} = 2\pi \frac{E}{\Delta E_{\text{cycle}}}$ & Measure of energy retention per cycle and frequency selectivity \\
RLC Resonance & $\omega_0 = 1/\sqrt{LC}$, $Q = \omega_0 L/R$ & Electrical analogue of the mechanical oscillator \\
\bottomrule
\end{tabularx}
\caption{Key relationships for damped and driven harmonic oscillators.}
\label{tab:ch09_summary}
\end{table}

\section{Chapter Exercises}
\begin{enumerate}
    \item \textbf{Underdamped Logarithmic Decrement:} Show that the ratio of two successive maximum displacements $x_n$ and $x_{n+1}$ of an underdamped oscillator separated by period $T = 2\pi/\omega_d$ satisfies $\delta_{\text{log}} \equiv \ln(x_n / x_{n+1}) = \gamma T = \frac{2\pi\gamma}{\sqrt{\omega_0^2 - \gamma^2}}$.
    \item \textbf{Bridge Resonance Safety:} An army detachment marches across a suspension bridge of mass $M = 50,000\text{ kg}$ and natural frequency $f_0 = 1.2\text{ Hz}$. Explain why troops are ordered to break step using the resonance amplitude formula.
    \item \textbf{Electrical-Mechanical Analogies:} Construct a dictionary table mapping mass $m$, damping $b$, spring $k$, position $x$, and driving force $F_0$ to their electrical series $LRC$ circuit counterparts.
    \item \textbf{Coupled Pendulums:} Two identical pendulums of length $\ell$ are connected by a spring of constant $k$ at their midpoints. Let $\theta_1$ and $\theta_2$ be the angular displacements. Write the linearised equations of motion, introduce normal coordinates $q_\pm = \theta_1 \pm \theta_2$, and identify the two normal mode frequencies $\omega_\pm$.
    \item \textbf{Critically Damped Shock Absorber:} A car's suspension system has spring constant $k = 12\;\text{kN/m}$ and effective mass $m = 300\;\text{kg}$. Calculate the critical damping coefficient $b_c = 2\sqrt{mk}$ needed to eliminate oscillation after a road bump, and write the general solution $x(t)$ for critical damping with initial conditions $x(0) = x_0$, $\dot{x}(0) = 0$.
\end{enumerate}
